\documentclass[11pt]{article}
\usepackage{amsmath, amssymb, amsthm}
\usepackage{enumerate}
\usepackage{mathpartir}
\usepackage{listings}
\usepackage{xcolor}
\usepackage{hyperref}

\definecolor{darkred}{HTML}{962d2d}
\definecolor{darkblue}{HTML}{2736b8}
\newcommand{\ctx}{\Gamma}
\newcommand{\synths}{\mathcolor{darkred}{\Rightarrow}}
\newcommand{\checks}{\mathcolor{darkblue}{\Leftarrow}}

\title{Type System Specification}
\author{Rem}

\begin{document}
\maketitle

\section{Simply Typed Lambda Calculus (STLC)}
We begin by formally defining the Simply Typed Lambda Calculus (STLC). This is the foundation of our
type system, which is extended with System F later on.

The syntax of types in STLC is nearly identical to the syntax of the lambda calculus itself.
The term $x$ is used to denote a variable, and likewise $x \vdash \tau$ denotes that the variable $x$
has type $\tau$ in the typing context $\ctx$.

\begin{mathpar}
    B = \{ \text{Int, Float, Bool, String, Unit, \dots} \} \\

    \tau ::= \tau \to \tau \mid T \quad \text{where } T \in B
\end{mathpar}

Let $B$ be the set of base types. The set of types $\tau$ is defined inductively over $B$, such that a type can either be a recursive type
definition of $\tau \to \tau$ or a base type $T$ where $T$ is an element of the set $B$.

The format for typing rules is as follows:
\begin{mathpar}
    \inferrule {premises}{conclusion} \quad \text{(Rule Name)}
\end{mathpar}

\subsubsection{Typing Rules for STLC}

\begin{mathpar}
    \inferrule {\mathit{x} {:} \sigma \in \ctx}{\ctx \vdash \mathit{x} {:} \sigma} \quad \text{(T-Var)}

    \inferrule {\text{$c$ is a constant of type $\ctx$}}{\ctx \vdash \mathit{c} {:} T}
    \quad \text{(T-Const)}

    \\
    \inferrule {\ctx, \mathit{x}{:}\sigma \vdash \mathit{e}{:}\tau}
    {\ctx \vdash (\lambda\mathit{x}{:}\sigma.\mathit{e}) {:} (\sigma \to \tau)} \quad \text{(T-Abs)}

    \inferrule {\ctx \vdash \mathit{e_1}{:} \sigma \to \tau \quad \ctx \vdash \mathit{e_2}{:}\sigma}
    {\ctx \vdash \mathit{e_1} \mathit{e_2}{:}\tau} \quad \text{(T-App)}
\end{mathpar}
\\
A typing assumption of $x:\sigma$ or $y:\tau$ means that the variable $x$ has either type $\sigma$ or $\tau$.

\pagebreak
In words, these typing rules can be described as follows:

\begin{itemize}
    \item T-Var: If $x$ has type $\sigma$ in the context $\ctx$, then $x$ has type $\sigma$
    \item T-Const: Term constants have their appropriate base types ($c \in B$)
    \item T-Abs: If under the assumption that $x$ has type $\sigma$, $e$ has type $\tau$, then in the same context without $x$, the lambda abstraction $\lambda x : \sigma . e$ has type $\sigma \to \tau$
    \item T-App: If $e_1$ has type $\sigma \to \tau$ and $e_2$ has type $\sigma$, then the application $e_1 e_2$ has type $\tau$
\end{itemize}

Examples of closed terms (terms typeable in the empty context) include:
\begin{itemize}
    \item $\lambda x : x.x$ (the identity function/I combinator)
    \item $\lambda x:\sigma.\lambda y:\tau.x$ (the K combinator)
    \item $\lambda t:\sigma.\lambda f:\sigma.t$ (Church encoding of the boolean value true)
\end{itemize}

\section{System F}
System F, also known as the polymorphic lambda calculus, extends STLC with universal quantification over types ($\forall$ quantification).
It was introduced by Jean-Yves Girard in 1972 and independently by John C. Reynolds in 1974.

\subsection{Typing Rules}
The typing rules for System F include all the rules from STLC, plus rules for type abstraction and type application.

\begin{mathpar}
    \inferrule {\ctx, \alpha \enspace \text{type} \vdash e {:} \tau}{\ctx \vdash \Lambda \alpha.e {:} \forall \alpha.\sigma}
    \quad \text{(T-TAbs)}

    \inferrule {\ctx \vdash e {:} \forall \alpha.\tau}{\ctx \vdash e \tau {:} \sigma[\tau/\alpha]}
    \quad \text{(T-TApp)}
\end{mathpar}

Where $\sigma,\tau$ are types, $\alpha$ is a type variable, and $\alpha \enspace \text{type}$ indicates that $\alpha$ is bound.
In words, these rules can be described as follows:
\begin{itemize}
    \item T-TApp: If $e$ has type $\forall \alpha.\tau$, then applying $e$ to a type $\tau$ results in a term of type $\sigma[\tau/\alpha]$, where $\sigma[\tau/\alpha]$ denotes the substitution of $\tau$ for $\alpha$ in $\sigma$.
    \item T-TAbs: If under the assumption that $\alpha$ is a type variable, $e$ has type $\tau$, then the type abstraction $\Lambda \alpha.e$ has type $\forall \alpha.\sigma$.
\end{itemize}

\subsubsection{Logic and Predicates in System F}
The Boolean type can be defined as $\forall \alpha.\alpha \to \alpha \to \alpha$ and we can implement Church encodings
for boolean and logical predicates as follows:

The functions for \textbf{True} and \textbf{False}, require three arguments --- a type $\alpha$ and two lambda expressions.
The two lambda expressions to decide which branch to take (either truthy or falsy).

\begin{enumerate}
    \item \textbf{T} = $\Lambda \alpha.\lambda x^\alpha.\lambda y^\alpha.x$
    \item \textbf{F} = $\Lambda \alpha.\lambda x^\alpha.\lambda y^\alpha.y$
\end{enumerate}

And then from these we can define the logical predicates as follows:
\begin{enumerate}
    \item \textbf{AND} = $\lambda x^{\text{Bool}}\lambda y^{\text{Bool}}.xy F$
    \item \textbf{OR} = $\lambda x^{\text{Bool}}\lambda y^{\text{Bool}}.xy T y$
    \item \textbf{NOT} = $\lambda x^{\text{Bool}}.x F T$
    \item \textbf{IF} = $\Lambda \alpha.\lambda x^{\text{Bool}}\lambda y^\alpha.\lambda z^\alpha.x \alpha yz$
    \item \textbf{ISZERO} = $\lambda n^{\text{Nat}}.n \text{Bool} (\lambda x^{\text{Bool}}.F) T$
\end{enumerate}

\subsubsection{Church Numerals}
System F allows recursive constructions to be embedded in a natural manner.
$K_1 \to K_2 \to \dots \to S$

Thus, we can define the type of natural numbers as follows using the Church encoding:
\begin{enumerate}
    \item \textit{zero} : $N$
    \item \textit{succ} : $N \to N$

    \item 0 := $\Lambda \alpha.\lambda x^\alpha.\lambda f^{\alpha \to \alpha}.x$
    \item 1 := $\Lambda \alpha.\lambda x^\alpha.\lambda f^{\alpha \to \alpha}.fx$
    \item 2 := $\Lambda \alpha.\lambda x^\alpha.\lambda f^{\alpha \to \alpha}.f(fx)$
    \item 3 := $\Lambda \alpha.\lambda x^\alpha.\lambda f^{\alpha \to \alpha}.f(f(fx))$
\end{enumerate}

A church numeral for $n$ is a higher-order function that takes a function $f$ and an argument $x$ and applies $f$ to $x$ $n$ times.

% source: https://arxiv.org/pdf/1908.05839
\section{Bidirectional Typing}
Bidirectional typing has two stages: type inference and type checking.
In our Haskell implementation these would be \texttt{infer} and \texttt{check} respectively.

\begin{itemize}
    \item Synthesis ($\synths$): The algorithm takes $\ctx$ and $e$ as input and produces a type $A$.
    \item Checking ($\checks$): The algorithm takes $\ctx$, $e$ and a known type $A$ as input and verifies if $e$ matches $A$.
\end{itemize}

\subsection{Subsumption}
Subsumption is the mechanism that allows the type system to switch from synthesis mode to checking mode.
It defines a rule where if an expression $e$ $\synths$ a type $A$, and $A = B$ then $e$ $\checks$ $B$.

\subsection{Bidirectional Typing Rules for STLC}

\begin{flalign*}
    \boxed{\ctx \vdash e {:} A}
    \quad \text{Under context $\ctx$, expression $e$ has type $A$}
\end{flalign*}

\begin{flalign*}
    \boxed{\ctx \vdash e \checks A}
     & \quad \text{Under $\ctx$, expression $e$ checks against type $A$} \\
    \boxed{\ctx \vdash e \synths A}
     & \quad \text{Under $\ctx$, expression $e$ synthesizes type $A$}
\end{flalign*}

\begin{mathpar}
    \inferrule {(x : A) \in \ctx}{\ctx \vdash x \synths A} \quad \text{(Var $\synths$)}
    \quad
    \inferrule {\ctx \vdash e \synths A \quad A = B}{\ctx \vdash e \checks B} \quad \text{(Subsumption $\checks$)}
\end{mathpar}

\begin{mathpar}
    \inferrule {\ctx \vdash e \checks A}{\ctx \vdash (e : A) \synths A} \quad \text{(Annot $\synths$)}
    \quad
    \inferrule {}{\ctx \vdash () \checks unit} \quad \text{(Unit $\checks$)}
    \\
    \inferrule {\ctx, x : A_1 \vdash e \checks A_2}{\ctx \vdash (\lambda x.e) \checks A_1 \to A_2} \quad \text{(LAbs $\checks$)}
    \\
    \inferrule {\ctx \vdash e_1 \synths A \to B \quad \ctx \vdash e_2 \checks A}
    {\ctx \vdash e_1 e_2 \synths B } \quad \text{(EApp $\synths$)}
\end{mathpar}

\pagebreak
In words,
\begin{itemize}
    \item Variables (Var $\synths$) synthesize their type from the context
    \item Subsumption (Subsumption $\checks$) allows switching from synthesis mode to checking mode
    \item Annotations (Annot $\synths$) allow switching from checking mode to synthesis mode
    \item Unit (Unit $\checks$) checks against the unit type '()'
    \item Lambda abstractions (LAbs $\checks$) are checked against function types by extending the context $\ctx$
          with the parameter type and checking the body against the return type
    \item Function applications (EApp $\synths$) synthesize their type by synthesizing the function type
          and checking the argument against the parameter type
\end{itemize}

\textbf{A note on lambda abstractions.}
Lambda abstractions (LAbs $\checks$) are checked against $A_1 \to A_2$ by \textit{extending} the context with $x : A_1$
and \textit{checking} the body $e$ against $A_2$. If we were to synthesize the type of a lambda abstraction,
we start with no information so we use a \textbf{checking rule} instead of a \textbf{synthesis rule}.

\subsection{Bidirectional Typing Rules for System F}
Extending the bidirectional rules from STLC to System F, requires adding two new rules to support universal
quantification over types ($\forall A.A$).

Type abstraction (T-TAbs) and type application (T-TApp).
Type abstraction is a checking rule, mirroing LAbs --- without a type annotation, there is no $\forall$ type to synthesize,
so we use a checking rule instead.

Type application is a synthesis rule, as the quantifier elimination allows us to synthesize the resulting type by substitution.


\begin{enumerate}
    \item EAbs Expr -- $\Lambda a. e$ (type abstraction, introduces $\forall$)
    \item ETyApp Expr Ty -- $e \enspace [A]$ (type application, eliminates $\forall$)
\end{enumerate}

\begin{mathpar}
    % T-TAbs checks (mirrors LAbs, forall introduction)
    \inferrule  {\ctx, \alpha \enspace \mathsf{type} \vdash e \checks \tau}
    {\ctx \vdash \Lambda \alpha.e \checks \forall \alpha.\tau}
    \quad \text{(T-TAbs $\checks$)}

    % T-TApp synthesizes (forall elimination)
    \inferrule {\ctx \vdash e \synths \forall\alpha.\tau}
    {\ctx \vdash e\enspace[A] \synths \tau[A/\alpha]}
    \quad \text{(T-TApp $\synths$)}
\end{mathpar}

\textbf{Comparison to the non-bidirectional typing rules.}
The non-bidirectional typing rules are only shown for comparison, as the bidirection rules
replace the previous rules.

\begin{mathpar}
    \inferrule {\ctx, \alpha \enspace \text{type} \vdash e {:} \tau}{\ctx \vdash \Lambda \alpha.e {:} \forall \alpha.\sigma}
    \quad \text{(T-TAbs)}

    \inferrule {\ctx \vdash e {:} \forall \alpha.\tau}{\ctx \vdash e \tau {:} \sigma[\tau/\alpha]}
    \quad \text{(T-TApp)}
\end{mathpar}

\textbf{A note on strong normalization ($\beta$-normal form).}
The typing rules do \textit{not} directly reduce terms to $\beta$-normal form, however strong normalization to $\beta$-normal form
is an inherent \textit{property} of the System F type system. We can conclude from this that all well-typed terms will \textit{eventually}
reduce to a $\beta$-normal form, but the typing rules themselves do not enforce this reduction.

% Bibtext references
\begin{thebibliography}{9}
    \bibitem{Girard's System F}
    \textit{Girard's System F} ---
    Jean-Yves Girard. (1972)

    \bibitem{Bidirectional Typing}
    \textit{Bidirectional Typing} ---
    Jana Dunfield and Neel Krishnaswami. (2020)
    \bibitem{Types and Programming Languages}
    \textit{Types and Programming Languages} ---
    Benjamin C. Pierce. (2002)

    \bibitem{STLC - Wikipedia}
    \textit{Simply Typed Lambda Calculus} ---
    \url{https://en.wikipedia.org/wiki/Simply_typed_lambda_calculus}

    \bibitem{System F - Wikipedia}
    \textit{System F} ---
    \url{https://en.wikipedia.org/wiki/System_F}
\end{thebibliography}

\end{document}